problem:  
Let \(n \ge 3\) and \(S^{n+1} \subset \mathbb{R}^{n+2}\) be the round sphere.  
Fix a compact group \( G \subset O(n+2)\) acting isometrically with cohomogeneity \(k \ge 2\), and let \(Q = S^{n+1}/G\) denote the quotient orbifold with its induced metric.  
Denote by \( \mathcal{M}_G(S^{n+1})\) the moduli space of embedded \(G\)-invariant minimal hyperspheres in \(S^{n+1}\), modulo congruence by ambient isometries.  
For each \(\Sigma \in \mathcal{M}_G(S^{n+1})\), let \(\operatorname{ind}_G(\Sigma)\) be the \(G\)-equivariant Morse index, i.e. the number of negative eigenvalues of the Jacobi operator restricted to \(G\)-invariant variations.

**Open Problem (Equivariant Index–Area Bounds):**  
Suppose \( \mathcal{M}_G(S^{n+1})\) contains a sequence of \(G\)-invariant embedded minimal hyperspheres \(\Sigma_j\) whose areas \(\mathcal{A}(\Sigma_j)\) tend to infinity.  
Does there exist a function \(F_G\colon [0,\infty)\to[0,\infty)\) and constants \(0<c_G<C_G<\infty\) (depending only on \(G,n\)) such that for all large \(j\),  
\[
c_G\,F_G(\mathcal{A}(\Sigma_j)) \le \operatorname{ind}_G(\Sigma_j) \le C_G\,F_G(\mathcal{A}(\Sigma_j))\,,
\]
and \(F_G(a)\) has a power-growth asymptotic comparable to \(a^{\beta_G}\) for some exponent \(\beta_G>0\)?  
More concretely: 
- Do such equivariant index–area inequalities exist, and 
- Is the growth rate \(\beta_G\) governed by analytic–geometric quantities of the orbit space \(Q\)—for example, by the asymptotics of the first \(G\)-invariant eigenvalues of the Laplace–Beltrami operator on \(Q\)?  

Why is it a "good" problem:  
This version corrects and focuses the conceptual core behind the earlier conjecture while removing vacuous or unjustified asymptotic assumptions.  
Rather than postulating a precise power law, it asks for **upper and lower comparison bounds** between the \(G\)-equivariant index and the area of symmetric minimal hyperspheres, parametrized by unknown growth \(F_G\).  

This question is:
- **Well posed analytically**, because the \(G\)-equivariant index is defined through the restriction of the stability operator to invariant variations, a natural elliptic operator tied to functions on the orbit space \(Q\).  
- **Realistic and open**, since even qualitative index–area comparability is unknown for symmetric minimal spheres, and such a bound would represent the first step toward a universal equivariant "stability law."  
- **Deeply interdisciplinary**, combining geometric measure theory, spectral analysis on singular quotient spaces, and the structure of \(G\)-invariant second variations.  

Discovering or disproving such index–area bounds would create a bridge between **spectral asymptotics** (through Laplace-like growth on \(Q\)) and **geometric stability** of minimal hypersurfaces under symmetry, illuminating how group actions quantitatively modify the landscape of instabilities.  
The statement remains conceptually simple (“does the equivariant index grow in controlled proportion to area under symmetry?”) while inviting deep analytic and geometric exploration—precisely the hallmarks of a *good* mathematical problem.